Both papers introduce a partial order $\preceq$ on
bin packing instances. Given two instances
$I=\{a_1,\ldots,a_n\}$ and $J=\{b_1,\ldots,b_m\}$, we write $I\preceq J$ if there exists an injective mapping
$
    f:\{1,\ldots,n\}\rightarrow\{1,\ldots,m\}
$
such that
$
    a_i\leq b_{f(i)},\,
    \forall i\in\{1,\ldots,n\}.
$
In other words, every item in $I$ can be matched
to a distinct item in $J$ of at least the same size.

Given the definition of the $\preceq$, we can obtain the \Cref{lem:monotonicity} which states that if $I \preceq J$ then states that $OPT(I) \leq OPT(J)$.


% \begin{lemma}[Monotonicity of Optimal Packing]
% \label{lem:monotonicity}
% Let $I=\{a_1,\ldots,a_n\}$ and
% $J=\{b_1,\ldots,b_m\}$ be two bin packing
% instances. If $I\preceq J$, then
% $\OPT(I)\leq\OPT(J).$
% \end{lemma}

% \begin{proof}
% By the definition of $I\preceq J$, there exists
% an injective mapping
% $f:\{1,\ldots,n\}\to\{1,\ldots,m\}$
% such that $a_i\leq b_{f(i)}$ for every
% $i\in\{1,\ldots,n\}$.

% Consider an optimal packing of $J$ using
% $\OPT(J)$ bins. For each item $a_i\in I$,
% replace the corresponding item $b_{f(i)}$
% with $a_i$, and remove all unmatched items
% from $J$.

% Since $a_i\leq b_{f(i)}$, the total size
% of items in each bin cannot increase.
% Hence, the resulting packing of $I$ is
% feasible and uses at most $\OPT(J)$ bins.

% Therefore,
% \[
%     \OPT(I)\leq\OPT(J).
% \]
% \end{proof}

